\subsection{Sprint 4}
Durante a Sprint 4, estavam previstas as seguintes \ac{us}: \ac{us}04 (Tela de Login), \ac{us}05 (Tela de Home), \ac{us}06 (Cadastro de Evento), \ac{us}09 (Cadastro de Usuário), \ac{us}10 (Cadastro de Roteiro), \ac{us}11 (Listagem de Roteiros), \ac{us}12 (Visualizar Roteiro), \ac{us}13 (Visualizar Mapa) e \ac{us}15 (Visualizar Perfil), além dos débitos técnicos remanescentes das sprints 2 e 3. Fiquei encarregada, junto com minha colega, da tarefa de visualização de rota turística.

A task da qual fiquei encarregada foi devidamente desenvolvida e concluída, assim como a maior parte das atividades previstas. A entrega da sprint foi realizada com sucesso e aceita, mas o time ainda está se dedicando para finalizar os últimos detalhes e apresentar um projeto que atenda plenamente às expectativas.

Durante essa sprint, enfrentamos desafios significativos. Devido à tragédia que ocorreu durante o andamento do projeto, um dos principais problemas que continuamos a enfrentar foi a falta de tempo. Muitas tarefas precisaram ser adiadas, resultando em um grande acúmulo de atividades para a última sprint e pouco tempo disponível para concluir todas as pendências. Essa situação nos ensinou importantes lições. Embora não possamos controlar os imprevistos que surgem, podemos controlar a maneira como lidamos com os obstáculos. Desse modo, demonstramos resiliência e comprometimento em entregar o melhor resultado possível.

Os últimos preparativos até a entrega final do projeto incluem finalizar alguns débitos técnicos e detalhes da documentação, além de preparar a apresentação e concluir os relatórios pendentes. Estamos determinados a superar os desafios e entregar um trabalho que reflita nosso esforço e dedicação. Também consideraremos o feedback recebido durante a revisão da sprint para planejar a implementação das melhorias finais. É importante que todas as sugestões dos \textit{stakeholders} sejam analisadas e, se possível, incorporadas ao projeto. A conclusão da documentação também será feita de forma detalhada, descrevendo cada etapa do processo de desenvolvimento e todas as instruções necessárias para a migração e continuidade do funcionamento da aplicação. Essa documentação será essencial para futuras referências e implementação do projeto. A equipe continua focada em entregar um projeto de alta qualidade, que não apenas atenda aos requisitos iniciais, mas que surpreenda positivamente. Acreditamos que, com a dedicação e o esforço coletivo, conseguiremos finalizar esta última etapa com sucesso e alcançar os objetivos propostos.